Ohne Wärmeverluste gibt das warme Wasser ($V_2$, $\vartheta_2$) so viel Wärme ab, wie das kalte Wasser ($V_1$, $\vartheta_1$) aufnimmt. Mit der Mischtemperatur $\vartheta$ gilt
\begin{align}
 c\,\rho V_2\left(\vartheta_2-\vartheta\right) &= c\,\rho V_1\left(\vartheta-\vartheta_1\right)
\end{align}
Spezifische Wärme und Dichte kürzen sich weg (der kleine Dichteunterschied zwischen warmem und kaltem Wasser ist vernachlässigt). Aufgelöst nach $\vartheta$:
\begin{align}
 \vartheta &= \TmF \\
   &= \frac{\VaO\cdot\TaO+\VbO\cdot\TbO}{\VaO+\VbO} \\
   &= \Tm \approx \result{\TmP{0}{2}}
\end{align}
Die Mischtemperatur ist das mit den Mengen gewichtete Mittel der beiden Temperaturen.
